\section{Ensemble Methods Overview}
\label{sec:ensemble-methods}

We review the three principal methods used to generate redistricting
ensembles, focusing on their outputs, limitations, and what each method
is suited to establish.

\subsection{GerryChain and the Flip Chain}

\gc{} \citep{gerrychain2021} implements a Markov chain on the space of
valid $k$-district plans for a given state graph.  The basic transition
is a \emph{flip}: a randomly chosen boundary node is reassigned to an
adjacent district, provided the resulting plan remains connected and within
population tolerance.  The chain is run for millions of steps; at each step
(or at a thinned subset of steps) the current plan is recorded.

\paragraph{Outputs.} The flip chain produces a sample from the
\emph{stationary distribution} of valid plans.  If run long enough and
sampled frequently enough, the histogram of any statistic (partisan seats,
Polsby-Popper score, minority district count) approximates the true
geometric distribution of that statistic over all valid plans.

\paragraph{Limitations.}
\begin{itemize}
  \item \textbf{Mixing time.} The flip chain mixes slowly on large state
    graphs (e.g., Texas at $k=38$).  Published ensemble sizes of 10,000
    to 1,000,000 plans may not achieve stationarity.
  \item \textbf{Locality bias.} Flips are local moves.  The chain may
    not explore topologically distinct map families within feasible runtime.
  \item \textbf{No canonical plan.} Any sampled plan is one draw from the
    distribution.  No plan in the ensemble has a privileged status.
\end{itemize}

\subsection{ReCom (Recombination)}

\recom{} \citep{deford2021recombination} replaces the flip transition with
a \emph{recombination} move: two adjacent districts are merged, then the
merged region is re-partitioned by a random spanning tree cut.  This
produces much larger structural changes per step.

\paragraph{Outputs.} \recom{} generates plans with different properties
than the flip chain.  In particular, \recom{} tends to produce more
``wiggly'' district boundaries, while the flip chain tends toward more
blocky shapes.  Both are valid samples from their respective stationary
distributions, but these distributions are not identical.

\paragraph{Limitations.}
\begin{itemize}
  \item \textbf{Non-reversibility.} Standard \recom{} is not reversible
    with respect to a known target distribution.  The stationary
    distribution is implicit.
  \item \textbf{Compactness underweighting.} The random spanning tree cut
    does not weight for compactness, so \recom{} ensembles often contain
    plans less compact than what a human mapmaker would draw.
  \item \textbf{Different distribution from flip chain.} Results from
    \recom{} and flip-chain ensembles cannot be directly compared without
    calibration.
\end{itemize}

The \citet{autry2021multiscale} Metropolized Multiscale Forest \recom{}
partially addresses the non-reversibility issue by applying a
Metropolis-Hastings correction, but at substantial computational cost.

\subsection{Short-Burst Methods}

Short bursts \citep{cannon2022voting} are a variant of \recom{} designed
specifically to optimise a target objective.  The chain is run for a
short fixed number of steps (a ``burst''), the best plan in the burst is
kept, and the process repeats.  Short bursts can hill-climb toward extreme
plans (e.g., maximum minority-district count) while remaining within the
convex hull of valid plans.

\paragraph{Outputs.} Short-burst ensembles are not samples from the
geographic distribution.  They are samples from the optimum-seeking
distribution around a specific objective.  They are useful for establishing
what is \emph{achievable} under the geographic constraints, not what is
\emph{typical}.

\paragraph{Limitations.}
\begin{itemize}
  \item Not a characterisation of the full plan space.
  \item Results depend strongly on the objective function and burst length.
  \item Cannot be used to place a plan in the ``neutral'' distribution.
\end{itemize}

\subsection{What Each Method Establishes}

Table~\ref{tab:method-comparison} summarises the three methods.

\begin{table}[h]
\centering
\caption{Comparison of ensemble methods and their evidentiary uses.}
\label{tab:method-comparison}
\begin{tabular}{lccc}
\toprule
Method & Stationary dist.\ known & Compactness weight & Best evidentiary use \\
\midrule
Flip chain       & Implicit (uniform) & None & Outlier detection \\
\recom{}         & Implicit           & None & Partisan distribution \\
Short-burst      & None (optimiser)   & None & Feasibility bounds \\
\ar{} (this work)& N/A (deterministic)& METIS edge cut & Canonical plan \\
\bottomrule
\end{tabular}
\end{table}

The \ar{} plan does not belong to any of these categories: it is not a
sample, but the output of a deterministic optimisation.  The G-track papers
use the \recom{} ensemble as the reference distribution and locate the
\ar{} plan within it.
